%==============================================================================
% Chapitre 19 - Organes de visée mécaniques
%==============================================================================

\section{Introduction}

Les organes de visée permettent au tireur d'orienter l'arme afin que le
projectile atteigne la cible visée.

Avant l'apparition des systèmes optiques modernes, cette fonction était
assurée exclusivement par des dispositifs mécaniques.

Aujourd'hui encore, les organes de visée mécaniques restent largement
utilisés :

\begin{itemize}
    \item comme système principal sur certaines armes ;
    \item comme système de secours ;
    \item pour l'instruction ;
    \item dans les applications où la robustesse est prioritaire.
\end{itemize}

\begin{aretenir}

La fonction première des organes de visée est de matérialiser une direction
de tir cohérente entre l'œil du tireur, l'arme et la cible.

\end{aretenir}

\section{Principe général}

Le projectile ne suit pas la ligne de visée.

En raison de la gravité, sa trajectoire est courbe.

Les organes de visée permettent donc de donner à l'arme une orientation telle
que le projectile rencontre la cible à une distance donnée.

Cette distance est appelée distance de réglage ou distance de zérotage.

\section{Ligne de visée}

La ligne de visée est la droite reliant :

\begin{itemize}
    \item l'œil du tireur ;
    \item les organes de visée ;
    \item la cible.
\end{itemize}

Dans la plupart des armes, cette ligne se situe au-dessus de l'axe du canon.

\begin{aretenir}

L'axe du canon et la ligne de visée ne sont généralement pas confondus.

\end{aretenir}

\section{Hausse et guidon}

Le système mécanique le plus répandu est constitué de :

\begin{itemize}
    \item la hausse ;
    \item le guidon.
\end{itemize}

Le guidon est situé à l'avant de l'arme.

La hausse est située à l'arrière.

Le tireur doit aligner :

\begin{enumerate}
    \item l'œil ;
    \item la hausse ;
    \item le guidon ;
    \item la cible.
\end{enumerate}

\section{Le guidon}

Le guidon constitue le repère avant du système de visée.

Il existe plusieurs formes :

\begin{itemize}
    \item lame ;
    \item pointe ;
    \item boule ;
    \item insert lumineux.
\end{itemize}

Son choix dépend du type d'arme et de l'emploi recherché.

\section{La hausse}

La hausse constitue le repère arrière.

Elle peut être :

\begin{itemize}
    \item fixe ;
    \item réglable ;
    \item à curseur ;
    \item à œilleton.
\end{itemize}

Certaines hausses permettent un réglage précis de l'élévation et de la
dérive.

\section{Visée ouverte}

Dans une visée ouverte classique :

\begin{itemize}
    \item le guidon apparaît centré dans l'encoche de hausse ;
    \item les sommets sont alignés ;
    \item la cible est observée au-delà du guidon.
\end{itemize}

Cette solution est simple et robuste mais nécessite un alignement précis.

\section{Visée à œilleton}

La visée à œilleton remplace l'encoche de hausse par une ouverture
circulaire.

L'œil a naturellement tendance à centrer le guidon dans cette ouverture.

Cette configuration offre souvent :

\begin{itemize}
    \item une meilleure précision ;
    \item une acquisition plus rapide ;
    \item une meilleure répétabilité.
\end{itemize}

\section{Réglage en site}

Le réglage en site correspond à l'ajustement vertical de la visée.

Il permet de compenser :

\begin{itemize}
    \item la chute du projectile ;
    \item la distance de tir ;
    \item certaines caractéristiques balistiques.
\end{itemize}

Traditionnellement, ce réglage est assuré par la hausse.

\section{Réglage en dérive}

Le réglage en dérive correspond à l'ajustement horizontal.

Il permet de corriger :

\begin{itemize}
    \item des défauts d'alignement ;
    \item certaines influences extérieures ;
    \item des écarts systématiques observés lors des tirs.
\end{itemize}

\section{Distance de zérotage}

Le zérotage consiste à régler les organes de visée afin que le point visé et
le point touché coïncident à une distance donnée.

Selon l'arme considérée, cette distance peut être :

\begin{itemize}
    \item 25 m ;
    \item 100 m ;
    \item 200 m ;
    \item 300 m ;
    \item ou toute autre valeur adaptée à l'usage prévu.
\end{itemize}

\section{Erreur de visée}

Les organes mécaniques introduisent plusieurs sources d'erreur :

\begin{itemize}
    \item mauvais alignement ;
    \item positionnement irrégulier de l'œil ;
    \item défaut de concentration ;
    \item mouvements du tireur.
\end{itemize}

Ces erreurs augmentent généralement avec la distance.

\section{Longueur de ligne de mire}

La distance séparant la hausse du guidon est appelée ligne de mire.

Une ligne de mire plus longue permet généralement :

\begin{itemize}
    \item une meilleure précision angulaire ;
    \item une réduction de certaines erreurs d'alignement.
\end{itemize}

C'est l'une des raisons pour lesquelles les armes longues disposent souvent
d'une précision de visée supérieure aux armes de poing.

\section{Influence sur la précision}

La précision finale dépend notamment :

\begin{itemize}
    \item de la qualité des organes de visée ;
    \item de leur réglage ;
    \item de la qualité du tireur ;
    \item des performances balistiques de l'arme.
\end{itemize}

Les organes de visée ne créent pas la précision intrinsèque de l'arme, mais
ils conditionnent la capacité du tireur à l'exploiter.

\section{Application aux simulations}

Dans les simulateurs de tir, les organes mécaniques sont souvent reproduits
avec un niveau de fidélité variable.

Les paramètres fréquemment représentés sont :

\begin{itemize}
    \item la géométrie du guidon ;
    \item la géométrie de la hausse ;
    \item les réglages ;
    \item les erreurs de pointage.
\end{itemize}

\begin{simulation}

Dans les simulateurs d'entraînement, la qualité de reproduction des organes
de visée influence directement le réalisme perçu par l'utilisateur.

\end{simulation}

\section{Vers les systèmes optiques}

Les organes mécaniques ont longtemps constitué la référence.

Les systèmes optiques modernes reprennent plusieurs notions introduites dans
ce chapitre :

\begin{itemize}
    \item ligne de visée ;
    \item réglage en site ;
    \item réglage en dérive ;
    \item distance de zérotage.
\end{itemize}

Ils apportent cependant des capacités supplémentaires qui seront étudiées
dans le chapitre suivant.

\section{À retenir}

\begin{aretenir}

\begin{itemize}
    \item Les organes de visée permettent d'orienter l'arme vers la cible.
    \item La ligne de visée est généralement distincte de l'axe du canon.
    \item La hausse et le guidon constituent le système mécanique le plus
    répandu.
    \item Le zérotage définit la distance pour laquelle la visée est réglée.
    \item Une ligne de mire longue améliore généralement la précision de
    visée.
\end{itemize}

\end{aretenir}
